\chapter{Prepare Python and the Project}

\begin{goalbox}
Create an isolated Python environment and install the one required package.
\end{goalbox}

Install Python 3.10 or newer. Check it in Terminal, PowerShell, or the terminal
inside your editor:

\begin{lstlisting}[language=bash,numbers=none]
python --version
\end{lstlisting}

On some macOS or Linux systems the command is \texttt{python3}. A result such as
\texttt{Python 3.12.4} is suitable.

\section{Open the code folder}

Clone the repository and enter its code folder:

\begin{lstlisting}[language=bash,numbers=none]
git clone https://github.com/alirezaaies/telbot-simple-bot.git
cd telbot-simple-bot/code
\end{lstlisting}

If you downloaded a ZIP file, extract it and open a terminal in its \texttt{code}
folder instead.

\section{Create an isolated environment}

An environment keeps this tutorial's package separate from other Python projects.

\textbf{macOS or Linux:}
\begin{lstlisting}[language=bash,numbers=none]
python3 -m venv .venv
source .venv/bin/activate
\end{lstlisting}

\textbf{Windows PowerShell:}
\begin{lstlisting}[language={},numbers=none]
python -m venv .venv
.venv\Scripts\Activate.ps1
\end{lstlisting}

The terminal usually adds \texttt{(.venv)} before its prompt. Then install the
dependency:

\begin{lstlisting}[language=bash,numbers=none]
python -m pip install --upgrade pip
python -m pip install -r requirements.txt
\end{lstlisting}

The requirements file contains only \texttt{pyTelegramBotAPI}. Its import name is
\texttt{telebot}; this difference is normal.

\begin{checkpointbox}
Run the command below. If it prints \texttt{telebot is ready}, installation is complete.
\begin{lstlisting}[language=bash,numbers=none]
python -c "import telebot; print('telebot is ready')"
\end{lstlisting}
\end{checkpointbox}
